\section{Evaluation}
\label{sec:evaluation}

\subsection{Models compared}

Four configurations are evaluated on a held-out test split derived from the synthetic
training data. They span two architectural choices --- a custom 125\,M-parameter
transformer versus the much larger pretrained Qwen2.5-3B --- and, within each choice, a
comparison of from-scratch (or zero-shot) against a model that has been adapted to the
task data. All four are wrapped behind the same retriever, so any difference in
generation quality is attributable to the model rather than to retrieval.

\begin{table}[h]
\centering
\small
\begin{tabular}{p{3.4cm}p{3.0cm}p{2.6cm}p{4.2cm}}
\toprule
\textbf{Model} & \textbf{Initialisation} & \textbf{Training data} & \textbf{Role} \\
\midrule
Custom (scratch) & Random ($\mathcal{N}(0, 0.02)$) & Synthetic + public QA &
Fulfils mandatory ``train one model from scratch'' criterion. \\
Custom (GPT-2 init) & GPT-2 small weights & Synthetic + public QA &
Same architecture, warm-started from a pretrained checkpoint. \\
Qwen2.5-3B-Instruct & Pretrained, frozen & --- & Strong zero-shot reference. \\
Qwen2.5-3B-Instruct + LoRA & Pretrained + LoRA adapters & Synthetic & Parameter-efficient
fine-tune on the synthetic task data \cite{hu2022lora}. \\
\bottomrule
\end{tabular}
\caption{The four configurations under comparison. Only the first two are decoder-only models
trained on this project's hardware; the second pair establishes how close a small custom model
gets to a much larger off-the-shelf instruction-tuned model.}
\label{tab:models}
\end{table}

The two custom configurations isolate the contribution of pretraining on a generic web
corpus, holding the architecture constant. The two Qwen configurations isolate the
contribution of LoRA fine-tuning on the synthetic data, holding the base model constant.

\subsection{Metrics}

Two families of metrics are reported: reference-based surface metrics and an LLM-as-judge
evaluation. They are complementary --- surface metrics are cheap and reproducible but
sensitive to surface form, while the judge captures qualitative properties such as
pedagogical clarity that no token-level metric can measure.

\paragraph{Surface metrics.}
ROUGE-L and Token-F1 were computed for completeness but are unreliable in this setting:
a concise reference and a verbose-but-correct prediction share few tokens, which
systematically penalises longer responses regardless of quality. BERTScore F1 is more
robust to paraphrase, but its range compresses near $0.85$ for all in-domain text and so
discriminates poorly between weak and strong outputs. METEOR is the most informative of
the surface metrics here because its stemming and synonym matching reduce the sensitivity
to length mismatch. Surface metrics are averaged over $n=200$ test examples.

\paragraph{LLM-as-judge.}
Each generated response is scored by \texttt{Qwen2.5-7B-Instruct} against a fixed rubric
on three axes --- \emph{correctness}, \emph{completeness}, and \emph{pedagogical
clarity} --- on an integer 1--5 scale. Judge scores are reported over $n=5$ examples per model;
the small sample size and the use of a Qwen-family judge alongside Qwen-family
configurations are discussed as limitations below.

\subsection{Results}

\begin{table}[h]
\centering
\caption{Evaluation results. Surface metrics computed over $n=200$ examples; judge
scores over $n=5$, graded by \texttt{Qwen2.5-7B-Instruct} on a 1--5 scale.}
\label{tab:eval}
\small
\setlength{\tabcolsep}{4pt}
\begin{tabular}{lccccc}
\toprule
\textbf{Model}
  & \textbf{METEOR} ($\uparrow$)
  & \textbf{BERTScore F1} ($\uparrow$)
  & \textbf{Correct.} ($\uparrow$)
  & \textbf{Complete.} ($\uparrow$)
  & \textbf{Clarity} ($\uparrow$) \\
\midrule
Custom (scratch)            & 0.157 & 0.855 & 1.4 & 1.2 & 1.2 \\
Custom (GPT-2 init)         & 0.208 & 0.862 & 2.4 & 1.4 & 1.2 \\
Qwen2.5-3B-Instruct         & 0.305 & 0.856 & 4.0 & 3.6 & \textbf{3.2} \\
Qwen2.5-3B-Instruct + LoRA  & \textbf{0.375} & \textbf{0.904} & \textbf{4.0} & \textbf{3.8} & 2.8 \\
\bottomrule
\end{tabular}
\end{table}

\subsection{Discussion}

The two 125\,M models sit at the bottom of every judge axis, confirming the suspicion
that surface-metric compression masks genuinely poor output: the BERTScore values for
all non-custom models fall within $0.855$--$0.904$, a spread that is still relatively
compressed compared to the much larger variation in judge scores. METEOR tracks the
judge ranking more faithfully and is the more informative of the two surface metrics in
this setting.

Initialising from GPT-2 weights nearly doubles the correctness score ($1.4 \to 2.4$),
consistent with a warm-started model that recovers relevant terminology faster, but
pedagogical clarity stays at $1.2$ in both small configurations. This is the qualitative
footprint of a model that knows the right words but cannot yet string them into a
coherent explanation --- a capacity ceiling rather than a training-time problem.

The Qwen2.5-3B baseline outperforms both 125\,M models on all judge axes by a wide
margin, confirming that on free-form correctness this project's custom model cannot
compete with a model trained on a web-scale corpus. LoRA fine-tuning on the synthetic
data improves both METEOR ($0.305 \to 0.375$) and BERTScore F1 ($0.856 \to 0.904$),
indicating stronger alignment with the reference answers.

Interestingly, the LoRA configuration matches the base model on correctness (4.0) and
slightly improves completeness ($3.6 \to 3.8$), but shows a drop in pedagogical clarity
($3.2 \to 2.8$). This suggests that while domain adaptation improves factual alignment
and coverage, it may bias the model toward more direct or terse answers at the expense
of explanation quality --- a common trade-off when fine-tuning on narrowly scoped
synthetic data.

Two methodological caveats apply. First, the judge sample size of $n=5$ is small;
the gap between the two custom models ($1.4$ vs.\ $2.4$ on correctness) and the gap
between the custom models and Qwen are large enough to survive that noise, but
finer-grained differences (such as the clarity drop under LoRA) should be read as
suggestive rather than statistically significant. Second, using a Qwen model as both
synthetic-data generator and as judge couples two otherwise independent steps in the
pipeline and may inflate the absolute scores of the Qwen-family configurations. A small
human evaluation, or a judge from a different model family (e.g.\ a Llama- or
Mistral-derived instruction model), would be the proper next step.
```
